% !PRINTED 68
\weberpage{68}
The totality of all rational functions of $r$ with rational numbers as coefficients is thus a number field $\Omega(r)$ of degree $\nu$, which we call a \gesperrt{cyclotomic field}. We wish, however, to take the concept of the cyclotomic field somewhat more generally, and to understand by it every field whose numbers are all rational functions of any roots of unity whatever.

The field $\Omega(r)$ of degree $\nu$, which consists of \gesperrt{all} rational functions of an $m$th root of unity $r$, we call, for more precise distinction, the \gesperrt{full cyclotomic field} of order $m$, and we denote it by $\Omega_m$.

Arbitrary roots of unity $r, r', r'' \ldots$ of arbitrary degrees $m, m', m'' \ldots$ can always be regarded as powers of one and the same root of unity $\varrho$, whose degree is the least common multiple of $m, m', m'' \ldots$. Accordingly a cyclotomic field which contains only rational functions of $r, r', r'' \ldots$ is a divisor of the full cyclotomic field $\Omega(\varrho)$, and we therefore obtain all cyclotomic fields that exist at all when we seek out all the divisors of all full cyclotomic fields.

The \gesperrt{Galois} group of the field $\Omega_m$ consists of all the substitutions
\weq{3}{(r, r^n),}
where $n$ denotes every relative prime to $m$, taken according to the modulus $m$. For the cyclotomic equation of degree $\nu$, whose roots are the $r^n$, is a normal equation and hence its own \gesperrt{Galois} resolvent. Since, if $a, b$ are two of these numbers $n$,
\[
(r, r^a)(r, r^b) = (r, r^{ab})
\]
this group is isomorphic with the group $\mathfrak{N}$ of all number classes $N$ of the numbers coprime to $m$, which we have considered in the preceding section.

If $\mathfrak{A}$ is a divisor of $\mathfrak{N}$, and $\varrho$ a function from $\Omega_m$ belonging to $\mathfrak{A}$, then the totality of the rational functions of $\varrho$, $\Omega(\varrho)$, is a field contained in $\Omega_m$.

If conversely $\Omega$ is a divisor of $\Omega_m$ and $\varrho$ a primitive number of the field $\Omega$, hence also a number in $\Omega_m$, then $\Omega$ can be represented as the totality of the rational functions of $\varrho$ and denoted by $\Omega(\varrho)$.
